% =====================================================================
% SECCIÓN 3: COMPRENSIÓN Y PREPARACIÓN DE LOS DATOS
% =====================================================================

% ---------------------------------------------------------------------
% Diapositiva 7: Corpus Textual Masivo
% ---------------------------------------------------------------------
\section{Comprensión y preparación de datos}

\begin{frame}{Corpus textual masivo: fuentes y consolidación}
  \framesubtitle{1\,311\,023 publicaciones de 260\,357 usuarios en español}

  \begin{table}
    \centering\scriptsize
    \setlength{\tabcolsep}{5pt}
    \caption{Composición del corpus unificado (\texttt{corpus\_v1.parquet}).}
    \begin{tabular}{lrrll}
      \toprule
      \textbf{Fuente} & \textbf{Documentos} & \textbf{\%} & \textbf{Idioma orig.} & \textbf{Nivel de etiqueta} \\
      \midrule
      Coello-Guilarte (Twitter) & 1\,047\,194 & 79{,}88\,\% & Español (nativo) & Usuario completo (0 / 2) \\
      Kaggle SDD (Reddit) & 232\,067 & 17{,}70\,\% & Inglés (trad.) & Comunidad / autoinforme \\
      Reddit Mental Health Posts & 24\,031 & 1{,}83\,\% & Inglés (trad.) & Subreddit temático \\
      Depression Reddit & 7\,731 & 0{,}59\,\% & Inglés (trad.) & Subreddit temático \\
      \midrule
      \textbf{Total consolidado} & \textbf{1\,311\,023} & \textbf{100\,\%} & -- & \textbf{260\,357 usuarios} \\
      \bottomrule
    \end{tabular}
  \end{table}

  \vskip-0.15cm
  \begin{columns}[T]
    \column{0.49\textwidth}
    \ugrblock{Anonimización y privacidad}{
      \begin{itemize}
        \item Supresión de URLs, menciones (@usuario), correos y teléfonos.
        \item Identificadores hasheados con SHA-256 (prefijo de 12 hex).
      \end{itemize}
    }
    \column{0.49\textwidth}
    \ugrblock{Calidad y multimodalidad}{
      \begin{itemize}
        \item Textos vacíos residuales inferiores al 0{,}02\,\%.
        \item Registro dual: microtexto de Twitter (14 tokens) vs.\ publicaciones extensas de Reddit (112 tokens).
      \end{itemize}
    }
  \end{columns}

  \note{
    Consolidamos un corpus inédito de más de 1,3 millones de textos anonimizados. La fuente nativa en español reúne 317 usuarios clínicos de Twitter, a los que sumamos tres colecciones de Reddit traducidas con redes neuronales para evaluar la transferencia de dominio. Toda la información personal fue eliminada de forma irreversible.
  }
\end{frame}

% ---------------------------------------------------------------------
% Diapositiva 8: Partición a Nivel de Usuario
% ---------------------------------------------------------------------
\begin{frame}{Partición estricta a nivel de usuario (prevención de fuga)}
  \framesubtitle{Estratificación 70/10/20 agrupada por autor (semilla fija 42)}

  \vskip-0.1cm
  \tikzdiagram{% Partición independiente por fuente: cada corpus se divide 70/10/20 por usuario.
\begin{tikzpicture}[font=\footnotesize, x=1cm, y=1cm]
    % Leyenda
    \node[anchor=west] at (0, 1.3) {%
        \tikz\fill[ugrAzul] (0,0) rectangle (0.3,0.3); Entrenamiento (70\,\%)\qquad
        \tikz\fill[ugrCeleste] (0,0) rectangle (0.3,0.3); Validación (10\,\%)\qquad
        \tikz\fill[ugrNaranja] (0,0) rectangle (0.3,0.3); Prueba (20\,\%)};

    % Barras por fuente
    \foreach \y/\nombre/\info in {
        0/Coello-Guilarte/317 usuarios,
        -1.05/Kaggle SDD/232\,067 usuarios sintéticos,
        -2.1/Reddit MH Posts/20\,242 usuarios,
        -3.15/Depression Reddit/7\,731 usuarios sintéticos} {
        \node[anchor=east, align=right] at (-0.25, \y) {\textbf{\nombre}\\[-1pt] {\color{ugrGris}\info}};
        \fill[ugrAzul]    (0, \y - 0.32) rectangle (7, \y + 0.32);
        \fill[ugrCeleste] (7, \y - 0.32) rectangle (8, \y + 0.32);
        \fill[ugrNaranja] (8, \y - 0.32) rectangle (10, \y + 0.32);
        \draw[white, line width=1.4pt] (7, \y - 0.32) -- (7, \y + 0.32) (8, \y - 0.32) -- (8, \y + 0.32);
        \node[text=white, font=\scriptsize\bfseries] at (3.5, \y) {usuarios completos, estratificados por clase};
    }

    % Test nativo común
    \draw[ugrRojo, line width=1.1pt, dashed, rounded corners=2pt] (7.92, -0.42) rectangle (10.08, 0.42);
    \node[ugrRojo, anchor=south, font=\scriptsize\bfseries] at (9.0, 0.44) {test nativo común};

    % Llaves
    \draw[decorate, decoration={brace, amplitude=5pt}, line width=0.8pt, ugrAzul] (10.35, 0.36) -- (10.35, -0.36)
        node[midway, right=7pt, text=ugrAzul] {\textbf{splits mono}};
    \draw[decorate, decoration={brace, amplitude=7pt}, line width=0.8pt, ugrNaranja!90!black] (12.75, 0.42) -- (12.75, -3.5)
        node[midway, right=9pt, align=left, text=ugrNaranja!90!black] {\textbf{splits multi}\\ unión de las\\ cuatro fuentes};
\end{tikzpicture}
}

  \vskip-0.15cm
  \begin{columns}[T]
    \column{0.49\textwidth}
    \ugralert{El riesgo de fuga (\emph{Data Leakage})}{
      Si se divide al azar por documento, textos del mismo autor caen en entrenamiento y prueba: el clasificador aprende a reconocer al autor en vez de la señal afectiva.
    }
    \column{0.49\textwidth}
    \ugrblockok{Blindaje metodológico}{
      Todos los textos de un usuario permanecen en un único conjunto. El \textbf{test nativo común} cuenta con \textbf{64 usuarios nunca vistos} en entrenamiento (250\,847 tweets reales).
    }
  \end{columns}

  \note{
    Uno de los pilares del trabajo es evitar el data leakage: agrupar la partición por autor. De esta forma, los 64 autores del conjunto de prueba nativo son completamente desconocidos para los modelos, garantizando una evaluación predictiva genuina sobre nuevos sujetos.
  }
\end{frame}

% ---------------------------------------------------------------------
% Diapositiva 9: Pipeline de Traducción y Preprocesamiento
% ---------------------------------------------------------------------
\begin{frame}{Pipeline de traducción y extracción de rasgos}
  \framesubtitle{Traducción neuronal y enriquecimiento psicolingüístico}

  \begin{columns}[T]
    \column{0.49\textwidth}
    \ugrblock{Traducción automática neuronal}{
      \begin{itemize}
        \item Modelo Transformer \textbf{MarianMT} (\texttt{opus-mt-en-es}).
        \item 263\,829 publicaciones traducidas con caché JSONL.
        \item Control de consistencia sintáctica y semántica.
      \end{itemize}
    }
    \vskip0.05cm
    \ugrblock{Balanceo de clases (solo entrenamiento)}{
      \begin{itemize}
        \item \textbf{Submuestreo 50/50:} 424\,436 textos en mono (usado en BETO para eficiencia en GPU).
        \item \textbf{Ponderación de clases:} 700\,717 textos en mono ($w_c \propto 1 / n_c$).
      \end{itemize}
    }

    \column{0.49\textwidth}
    \ugrblock{19 Rasgos psicolingüísticos y temporales}{
      Inspirados en la literatura psiquiátrica (LIWC) adaptados al español:
      \begin{itemize}
        \item \textbf{Foco en sí mismo:} Pronombres y verbos en 1ª persona del singular (\emph{yo, me, mi, conmigo}).
        \item \textbf{Pensamiento absolutista:} Términos cognitivos rígidos (\emph{siempre, nunca, jamás, nada, todo}).
        \item \textbf{Léxico afectivo:} Emociones NRC en español y polaridad léxica VADER.
        \item \textbf{Patrones temporales:} Hora de publicación y signos de insomnio o vigilia nocturna.
      \end{itemize}
    }
  \end{columns}

  \note{
    Para traducir Reddit utilizamos MarianMT en un pipeline local con caché para asegurar reproducibilidad. Paralelamente, implementamos un módulo para extraer 19 rasgos psicolingüísticos clásicos basados en Pennebaker: autorreferencia excesiva, absolutismo y polaridad afectiva, excluyendo cualquier rasgo derivado de la cantidad de mensajes para evitar fuga.
  }
\end{frame}
